Partial Differential Equations (PDEs) provide fundamental models of
physical phenomena evolving in space and time. Seminal examples include
Fourier's theory of heat conduction~\cite{fourier1822chaleur},
d'Alembert's wave equation for vibrating strings~\cite{dalembert1747vibration},
and Stokes' equations for viscous fluid motion~\cite{stokes1880friction}.
These models form a basis for understanding thermal, mechanical, and
fluid processes across science and engineering.
Controlling these processes requires stability and performance guarantees
that account for uncertain parameters, changing operating conditions, and
nonlinear effects.
Establishing these guarantees is challenging because both the dynamics
and the uncertain or nonlinear interactions may be distributed over space.

A systematic treatment of these distributed dynamics begins with
infinite-dimensional systems theory, which provides foundations for
studying PDE solutions, input--output maps, and well-posed feedback
interconnections~\cite{curtain1978infinite,curtain1995introduction,staffans2005wellposed}.
Lasiecka and Triggiani~\cite{lasiecka2000parabolic} developed continuous
and approximation theories for optimal control of abstract parabolic
systems.
Developments in robust stabilization and $H_\infty$
control~\cite{curtain1986robust,vankeulen1993hinfinity} are complemented by
backstepping and Lyapunov-based
designs~\cite{krstic2008boundary,fridman2012sampled}, and semidefinite
methods for PDE stability analysis~\cite{valmorbida2016pdesdp}.
A complementary computational route uses established robust-control tools
for Ordinary Differential Equations (ODEs), including frequency-domain
methods, Riccati equations, and Linear Matrix Inequalities
(LMIs)~\cite{Francis1987Course,doyle1989statespace,boyd1994lmi,zhou1996robust,packard1993complex}.
These tools apply to PDEs through spatial discretization, modal
truncation, or lumped-parameter approximation. Although convergence and
performance guarantees are available under suitable
conditions~\cite{morris2001hinfinityapproximation}, neglected modes can
affect the closed loop through spillover~\cite{balas1978modal}. This
motivates computational methods that retain the distributed dynamics
when constructing control certificates.

To retain distributed dynamics while exploiting such computational tools,
the Partial Integral Equation (PIE)
representation~\cite{peet2021pierepresentation,shivakumar2024GDPE}
generalizes ODE state-space models to coupled finite-dimensional and distributed
states. For supported PDE classes, it exactly incorporates the dynamics
and boundary conditions into bounded Partial Integral (PI) operators.
Their algebra is closed under addition, composition, and adjoints,
enabling extensions of matrix-based ODE methods. Linear PI Inequalities (LPIs) generalize LMIs, and
PIETOOLS~\cite{9147712,pietools2025manual} converts parameterized LPI
conditions into semidefinite programs while retaining the PIE plant
representation. An industrial application is the PIE based observer design for
the digital twin
of an inkjet printer's thermo-fluidic fixation process developed and
validated by Toolhally et al.~\cite{TOOLHALLY2026103558}.
Recently, Shivakumar et al.~\cite{shivaknew} established PIE duality and
convex LPI conditions for nominal $H_\infty$-optimal state-feedback
synthesis. This motivates extending PIE duality to uncertainty and
nonlinearities described by structural information beyond unstructured
small-gain conditions.

This approach builds on Lur'e and Postnikov's absolute-stability
problem~\cite{lure1944stability}, Popov's circle
criterion~\cite{popov1961absolute}, and Zames's input--output stability
theory~\cite{zames1966input}.
Zames--Falb multipliers~\cite{zames1968stability} exploit monotonicity and
slope restrictions.
Megretski and Rantzer~\cite{megretski1997system} established the modern
Integral Quadratic Constraint (IQC) framework, unifying such descriptions
for analysis of linear systems interconnected with uncertain or nonlinear
components. Within the PIE framework, hard IQCs
provide robust stability and performance tests for coupled ODE--PDE
systems, including infinite-dimensional uncertainty
channels~\cite{9303892,10156335,callegari}. Subsequent results connect PIE IQCs to
structured singular-value bounds and robust observer with known uncertainty
synthesis~\cite{lenssen2025muanalysissynthesisframeworkpartial}.
We develop a dual IQC representation for these analysis certificates and
apply it to convex state-feedback synthesis.

We express them in the
language of Willems's dissipation theory~\cite{willems1972dissipative1}.
For finite-dimensional linear systems, the famous Kalman--Yakubovich--Popov
(KYP) lemma~\cite{rantzer1996kyp} connects frequency-domain inequalities
to quadratic dissipation inequalities expressed as LMIs.
Hard factorizations connect IQCs to dissipativity by expressing them as
integral inequalities valid over every finite time horizon. This allows
IQC analysis to be formulated through dissipation inequalities.
Goh~\cite{577607} linked IQC to generalized sectors,
dissipativity and $J$-spectral factorization.
For rational IQC multipliers, existence conditions for $J$-spectral
factors are given by the canonical factorization theorem of
Bart, Gohberg, and Kaashoek~\cite{Bart1979MinimalFactorization} and by
Meinsma's characterization using equalizing vectors~\cite{MEINSMA1995243}.
Seiler~\cite{6915700} used the latter to establish existence for strict
positive--negative multipliers and showed that these factorizations yield
hard IQCs and dissipation inequalities with nonnegative storage,
related developments include doubly-hard
factorizations and terminal-cost
IQCs~\cite{carrasco2018equivalence,scherer2018dynamicdissipation}.
For fixed filters and controllers, these certificates yield convex
analysis conditions.

Designing the state-feedback gain introduces a synthesis obstacle already
visible for ODEs. For a fixed IQC filter, consider the augmented dynamics
\[
    \dot x=Ax+Bw+B_\mrm{u}u,\qquad
    v=Cx+Dw+D_\mrm{zu}u,
\]
where $v$ enters the IQC. Under state feedback $u=Kx$, a sufficient
dissipation condition is
\[
\begin{gathered}
    \begin{bmatrix}
        \bigl(P(A+B_\mrm{u}K)\bigr) +  \bigl(P(A+B_\mrm{u}K)\bigr)^\top& PB\\
        B^\top P & \epsilon I
    \end{bmatrix}\\
    {}+
    \begin{bmatrix}
        (C+D_\mrm{zu}K)^\top\\ D^\top
    \end{bmatrix}
    V
    \begin{bmatrix}
        C+D_\mrm{zu}K & D
    \end{bmatrix}
    \preceq0,
\end{gathered}
\]
where $P=P^\top\succeq0$ is the storage matrix, $V=V^\top$ is the
multiplier weight, and $\epsilon>0$.
For fixed $K$, this is an LMI. Optimizing $K$ couples it to both $P$ and
$V$ through $PB_\mrm{u}K$ and $VD_\mrm{zu}K$.
The standard ODE invertible congruence transformation removes the first coupling but
generally leaves the second. Existing methods use additional structure
or iterative design~\cite{7353149}, dual IQCs also permit convex
estimator and feedforward synthesis~\cite{kose2009robust,venkataraman} through the 
elimination lemma \cite{Meijer_2024}.

Our main result is a dual IQC construction through $J$-spectral
factorization, combined with PIE duality. The construction supports
infinite-dimensional channels and memoryless or dynamic filters.
As an application, we derive a convex LPI condition for robust
state-feedback synthesis with explicit controller recovery. The resulting
controller guarantees stability and the prescribed quadratic performance
for all uncertainties satisfying the IQC.

The paper is organized as follows.
Sections~\ref{sec:preliminaries},\ref{sec:problem-definition} introduce the
preliminaries, problem definition and standing assumptions. Section~\ref{sec:iqc-theory}
establishes robust stability and performance through the IQC framework.
Section~\ref{sec:primal-controller-analysis} examines the primal
analysis condition, and
Sections~\ref{sec:duality}--\ref{sec:dual-controller-synthesis}
develop duality and convex controller synthesis.
Sections~\ref{sec:numerical-examples}
and~\ref{sec:finite-dimensional-multipliers} present numerical examples
and connections to finite-dimensional dynamic multipliers.
Section~\ref{sec:conclusion} concludes the paper, proofs are collected
in the appendix.
